
\paragraph{Extended Gauge Framework (V1)}

\[
\bm{u}(x) = x(a_\varepsilon + \eta_\varepsilon)\FrameBasis + x\GaugeShearB - \tfrac{1}{2}x^2\bm{a}_{\Section} - \tfrac{1}{6}x^3\bm{b}_{\Section} - \tfrac{1}{2}x^2(\FrameBasis\otimes\CentroidLocation)\bm{b}_{\Section}
\]

\[
\bm{\theta}(x) = \eta_\vartheta\FrameBasis - \FrameBasis\times\GaugeShearA + x(a_\vartheta + c_\vartheta)\FrameBasis - x(\FrameBasis\times\bm{a}_{\Section}) - \tfrac{1}{2}x^2(\FrameBasis\times\bm{b}_{\Section})
\]


\paragraph{Extended Kinematic Fields (V2)}

To achieve section-averaged kinematics, we introduce an additional gauge
parameter \(\GaugeShearB \in \mathbb{R}^2\) that modifies the transverse
velocity:
\[
\bm{u}(x) = x(a_\varepsilon + \eta_\varepsilon)\FrameBasis + x\GaugeShearB - \tfrac{1}{2}x^2\bm{a}_{\Section} - \tfrac{1}{6}x^3\bm{b}_{\Section} - \tfrac{1}{2}x^2(\FrameBasis\otimes\CentroidLocation)\bm{b}_{\Section}\]
\[\bm{\theta}(x) = \eta_\vartheta\FrameBasis - \FrameBasis\times\GaugeShearA + x(a_\vartheta + c_\vartheta)\FrameBasis - x(\FrameBasis\times\bm{a}_{\Section}) - \tfrac{1}{2} x^2(\FrameBasis\times\bm{b}_{\Section})
\]

The compensating warping modes are:

\begin{longtable}[]{@{}
  >{\raggedright\arraybackslash}p{(\linewidth - 6\tabcolsep) * \real{0.1250}}
  >{\raggedright\arraybackslash}p{(\linewidth - 6\tabcolsep) * \real{0.3958}}
  >{\raggedright\arraybackslash}p{(\linewidth - 6\tabcolsep) * \real{0.2292}}
  >{\raggedright\arraybackslash}p{(\linewidth - 6\tabcolsep) * \real{0.2500}}@{}}
\toprule\noalign{}
\begin{minipage}[b]{\linewidth}\raggedright
Mode
\end{minipage} & \begin{minipage}[b]{\linewidth}\raggedright
Shape \(\bm{\Phi}\)
\end{minipage} & \begin{minipage}[b]{\linewidth}\raggedright
Amplitude
\end{minipage} & \begin{minipage}[b]{\linewidth}\raggedright
Constraint
\end{minipage} \\
\midrule\noalign{}
\endhead
\bottomrule\noalign{}
\endlastfoot
\(r\) & \(\FrameBasis\otimes\bm{r}\) &
\(\bm{\alpha}_r = -\GaugeShearA\) &
\(\mathbf{L}_{\rm vn}^{(r)}\bm{\gamma}\) \\
\(v\) & \(\oneS\) & \(\bm{\alpha}'_v = -\GaugeShearB\) &
\(\mathbf{L}_{\rm wn}^{(v)}\bm{\gamma}\) \\
\end{longtable}

\begin{longtable}[]{@{}
  >{\raggedright\arraybackslash}p{(\linewidth - 6\tabcolsep) * \real{0.1132}}
  >{\raggedright\arraybackslash}p{(\linewidth - 6\tabcolsep) * \real{0.3585}}
  >{\raggedright\arraybackslash}p{(\linewidth - 6\tabcolsep) * \real{0.2075}}
  >{\raggedright\arraybackslash}p{(\linewidth - 6\tabcolsep) * \real{0.3208}}@{}}
\toprule\noalign{}
\begin{minipage}[b]{\linewidth}\raggedright
Mode
\end{minipage} & \begin{minipage}[b]{\linewidth}\raggedright
Shape \(\bm{\Phi}\)
\end{minipage} & \begin{minipage}[b]{\linewidth}\raggedright
Amplitude
\end{minipage} & \begin{minipage}[b]{\linewidth}\raggedright
Constraint Type
\end{minipage} \\
\midrule\noalign{}
\endhead
\bottomrule\noalign{}
\endlastfoot
\(r\) & \(\FrameBasis\otimes\bm{r}\) &
\(\bm{\alpha}_r = -\GaugeShearA\) & Algebraic \\
\(v\) & \(\mathbf{1}_{\Section}\) & \(\bm{\alpha}_v = -x\GaugeShearB\) &
Differential \\
\end{longtable}

\subsubsection{Resultants}
\paragraph{Mode v (V1)}

Shape: \(\bm{\Phi}_v = \oneS\) (constant transverse)

This is purely transverse, so \(\FrameBasis\cdot\bm{\Phi}_v = 0\), hence
\(\omega_v = 0\) and \(\nabla\omega_v = 0\).

Strain contribution from this mode:
\[\hat{\bm{\varepsilon}}_{(v)} = \bm{\Phi}_v\bm{\alpha}'_v = \oneS\bm{\alpha}'_v = -\GaugeShearB\]

This is a constant transverse addition to the shear strain field.

\textbf{Bimoment} (conjugate to \(\bm{\alpha}'_v\)):
\[\bm{w}_{(v)} = \int_{\Section} \bm{\sigma} \cdot \frac{\partial\hat{\bm{\varepsilon}}}{\partial\bm{\alpha}'_v} \, dA = \int_{\Section} \bm{\sigma} \cdot \oneS \, dA = \int_{\Section} \bm{\tau} \, dA = \ShearForce\]

\textbf{Bishear} (conjugate to \(\bm{\alpha}_v\)):

The strain doesn't depend on \(\bm{\alpha}_v\) directly (only on
\(\bm{\alpha}'_v\)), and \(\nabla\omega_v = 0\). So:
\[
\bm{v}_{(v)} = 0
\]

Equilibrium is verified:
\[
\bm{w}'_{(v)} - \bm{v}_{(v)} = \ShearForce' - {0} = {0}
\]

(Under Saint-Venant loading with no distributed transverse force)
\paragraph{Warping Resultants (V2)}

\textbf{Mode \(v\):} The shape \(\bm{\Phi}_v = \mathbf{1}_{\Section}\) is
purely transverse with \(\bm{\Phi}_{v,\parallel} = 0\).

The bimoment, conjugate to \(\bm{\alpha}'_v\), is:
\[
\bm{w}_{(v)} = \int_{\Section} \bm{\tau}\,dA = \ShearForce
\]

The bishear, conjugate to \(\bm{\alpha}_v\), vanishes since
\(\nabla\bm{\Phi}_{v,\parallel} = 0\): 
\[
\bm{v}_{(v)} = 0
\]

Equilibrium under Saint-Venant loading:
\(\bm{w}'_{(v)} - \bm{v}_{(v)} = \ShearForce' = \mathbf{0}\).

\textbf{Mode \(r\):} The shape \(\bm{\Phi}_r = \FrameBasis\otimes\bm{r}\) 
has \(\bm{\Phi}_{r,\parallel} = \bm{r}^T\) and
\(\nabla\bm{\Phi}_{r,\parallel} = \mathbf{1}_{\Section}\).

The bimoment is:
\[
\bm{w}_{(r)} = \int_{\Section} \sigma_{ax}\,\bm{r}\,dA = Ex\,\bm{I}_m\bm{b}_{\Section}
\]

The bishear is: 
\[
\bm{v}_{(r)} = \int_{\Section} \bm{\tau}\,dA = \ShearForce
\]

Equilibrium:
\(\bm{w}'_{(r)} - \bm{v}_{(r)} = E\bm{I}_m\bm{b}_{\Section} - \ShearForce = 0\),
using the Saint-Venant relation \(\ShearForce = E\bm{I}_m\bm{b}_{\Section}\).

\subsubsection{V1}

\paragraph{Cowper Identification (V1)}


From
\(\bar{\bm{u}}'_\perp = (-\tfrac{1}{2}x^2 + \SectionAverage{\PoissonFlexure})\bm{b}_{\Section}\)
and \(\bm{u}'_\perp = \GaugeShearB - \tfrac{1}{2}x^2\bm{b}_{\Section}\):
\[
\GaugeShearB = \SectionAverage{\PoissonFlexure}\bm{b}_{\Section}
\]

From
\(\bar{\bm{\theta}}_\perp = \bar{\Theta}[\FrameBasis\otimes\bm{\psi}]\bm{b}_{\Section} - \tfrac{1}{2}x^2(\FrameBasis\times\bm{b}_{\Section})\)
and
\(\bm{\theta}_\perp = -\FrameBasis\times\GaugeShearA - \tfrac{1}{2}x^2(\FrameBasis\times\bm{b}_{\Section})\):
\[
-\FrameBasis\times\GaugeShearA = \bar{\Theta}[\FrameBasis\otimes\bm{\psi}]\bm{b}_{\Section}
\]
\[
\GaugeShearA = \FrameBasis\times\bar{\Theta}[\FrameBasis\otimes\bm{\psi}]\bm{b}_{\Section}
\]

Then:
\[
\bm{\gamma} = \GaugeShearB + \GaugeShearA = \SectionAverage{\PoissonFlexure}\bm{b}_{\Section} + \FrameBasis\times\bar{\Theta}[\FrameBasis\otimes\bm{\psi}]\bm{b}_{\Section} = (\ScaleShearGeomA)^{-1}\bm{b}_{\Section} = \bar{\bm{\gamma}}
\]

This confirms
\((\ScaleShearGeomA)^{-1} = \SectionAverage{\PoissonFlexure} + \FrameBasis\times\bar{\Theta}[\FrameBasis\otimes\bm{\psi}]\).

Define: \[
\ScaleShearGeomB := -\SectionAverage{\PoissonFlexure}\ScaleShearGeomA\]
\[\bar{\Xi}_\theta := (\FrameBasis\times\bar{\Theta}[\FrameBasis\otimes\bm{\psi}])^{-1}
\]

So \(\GaugeShearB = -(\ScaleShearGeomB)^{-1}\bm{\gamma}\) and
\(\GaugeShearA = \bar{\Xi}_\theta^{-1}\ScaleShearGeomA\bm{\gamma}\).

Note:
\((\ScaleShearGeomB)^{-1} + \bar{\Xi}_\theta^{-1}\ScaleShearGeomA = (\ScaleShearGeomA)^{-1}\),
i.e.,
\(-\SectionAverage{\PoissonFlexure} + \FrameBasis\times\bar{\Theta}[\FrameBasis\otimes\bm{\psi}]\ScaleShearGeomA = (\ScaleShearGeomA)^{-1}\)\ldots{}

Wait, let me verify:
\(\GaugeShearB + \GaugeShearA = -(\ScaleShearGeomB)^{-1}\bm{\gamma} + \bar{\Xi}_\theta^{-1}\ScaleShearGeomA\bm{\gamma}\)

We need this to equal \(\bm{\gamma}\). 
With
\(\bm{\gamma} = (\ScaleShearGeomA)^{-1}\bm{b}_{\Section}\):

\[
\GaugeShearB + \GaugeShearA = \SectionAverage{\PoissonFlexure}\bm{b}_{\Section} + \FrameBasis\times\bar{\Theta}[\FrameBasis\otimes\bm{\psi}]\bm{b}_{\Section} = (\ScaleShearGeomA)^{-1}\bm{b}_{\Section} = \bm{\gamma}
\]




\subsubsection{V2}

\paragraph{Identification (V2)}

Matching the transverse deflection \(\bm{u}_\perp = \bar{\bm{u}}_\perp\)
requires: 
\[
\GaugeShearB = \SectionAverage{\PoissonFlexure}\bm{b}_{\Section}
\]

Matching the transverse rotation
\(\bm{\theta}_\perp = \bar{\bm{\theta}}_\perp\) requires:
\[
\GaugeShearA = \FrameBasis\times\bar{\Theta}[\FrameBasis\otimes\bm{\psi}]\bm{b}_{\Section}
\]

These identifications yield the section-averaged shear strain:
\[
\bm{\gamma} = \left(\SectionAverage{\PoissonFlexure} + \FrameBasis\times\bar{\Theta}[\FrameBasis\otimes\bm{\psi}]\right) \bm{b}_{\Section}
= (\ScaleShearGeomA)^{-1}\bm{b}_{\Section}
\]

Inverting furnishes the constraint operators in terms of
\(\bm{\gamma}\):
\[
\GaugeShearB = \SectionAverage{\PoissonFlexure}\ScaleShearGeomA\bm{\gamma}, \qquad \GaugeShearA = \FrameBasis\times\bar{\Theta}[\FrameBasis\otimes\bm{\psi}]\ScaleShearGeomA\bm{\gamma}
\]

\paragraph{Constraints Matrices (V1)}

For Cowper, the gauge parameters are:
\[
\begin{aligned}
\GaugeShearB &= \SectionAverage{\PoissonFlexure}\ScaleShearGeomA\bm{\gamma}\\
\GaugeShearA &= \FrameBasis\times\bar{\Theta}[\FrameBasis\otimes\bm{\psi}]\ScaleShearGeomA\bm{\gamma}
\end{aligned}
\]

The constraint operators \(\bm{\alpha}'_v = -\GaugeShearB\) and 
\(\bm{\alpha}_r = -\GaugeShearA\) are
\[
\begin{aligned}
\mathbf{L}_{\rm wn}^{(v)} &= -\SectionAverage{\PoissonFlexure}\ScaleShearGeomA \\
\mathbf{L}_{\rm vn}^{(v)} &= \mathbf{0}
\end{aligned}
\qquad\text{and}\qquad
\begin{aligned}
\mathbf{L}_{\rm vn}^{(r)} &= -\FrameBasis\times\bar{\Theta}[\FrameBasis\otimes\bm{\psi}]\ScaleShearGeomA \\
\mathbf{L}_{\rm wn}^{(r)} &= \mathbf{0}
\end{aligned}
\]

\paragraph{Constraint Operators (V2)}
The \(v\) and \(r\) amplitues are differentially and algebraically constraind, respectively:
\[
\begin{aligned}
\mathbf{L}_{\rm wn}^{(v)} &= -\SectionAverage{\PoissonFlexure}\ScaleShearGeomA, && \mathbf{L}_{\rm vn}^{(v)} = 0 \\
\mathbf{L}_{\rm vn}^{(r)} &= -\FrameBasis\times\bar{\Theta}[\FrameBasis\otimes\bm{\psi}]\ScaleShearGeomA, 
&& \mathbf{L}_{\rm wn}^{(r)} = 0
\end{aligned}
\]

\paragraph{Axial Rotation (V2)}

The axial component of rotation in the gauged model is
\(\vartheta = \eta_\vartheta + x(a_\vartheta + c_\vartheta)\), while
Cowper's averaged axial rotation is 
\(\bar{\vartheta} = x\,\FrameBasis\cdot\bar{\Theta}[\PoissonFlexure]\bm{b}_{\Section}\).

Under pure flexure, matching would require
\(c_\vartheta = \FrameBasis\cdot\bar{\Theta}[\PoissonFlexure]\bm{b}_{\Section}\).
However, the Saint-Venant solution determines:
\[
\Jsv c_\vartheta = -\FrameBasis\cdot\int_{\Section} \bm{r}\times(\nabla\bm{\psi} + \PoissonFlexure)\bm{b}_{\Section} \dA
\]

These expressions differ by the shear warping contribution
\(\nabla\bm{\psi}\), which represents shear-induced torsion in
asymmetric sections. 
To maintain Saint-Venant exactness, \(\vartheta\)
should retain this coupling rather than match the section average.


\subsection{Axial Rotation}

\paragraph{Resultants}

With \(\bm{\Phi}_{\vartheta,\parallel} := \FrameBasis\cdot(\FrameBasis\times\bm{r}) = 0\),
the bimoment conjugate to \(\alpha'_\vartheta\) is:
\[
w_{(\vartheta)} := \int_{\Section} \bm{\tau}\cdot(\FrameBasis\times\bm{r})\,dA = \FrameBasis\cdot\int_{\Section} \bm{r}\times\bm{\tau}\,dA = T
\]
ie, the torque resultant.

\textbf{Bishear} (conjugate to \(\alpha_\vartheta\)):

\(\nabla\bm{\Phi}_{\vartheta,\parallel} = 0\), so:
\[
v_{(\vartheta)} = 0
\]

\textbf{Equilibrium:} 
\[
w'_{(\vartheta)} - v_{(\vartheta)} = T' = 0
\]
under SV.

The identification would be:
\[\tilde{\eta}_\vartheta = c_\vartheta - \bar{c}_{\vartheta}\]

where:
\[\bar{c}_{\vartheta} := \FrameBasis\cdot\bar{\Theta}[\PoissonFlexure]\bm{b}_{\Section} = \FrameBasis\cdot\bar{\bm{I}}^{-1}\int_{\Section} (\bm{r}-\bar{\bm{p}})\times\PoissonFlexure\,dA\,\bm{b}_{\Section}\]

We need to express this in terms of \(\bm{\gamma}\). Using
\(\bm{b}_{\Section} = \bar{\Xi}_\kappa\bm{\gamma}\):
\[\tilde{\eta}_\vartheta = \left(c_\vartheta/\bm{b}_{\Section} - \bar{c}_{\vartheta}/\bm{b}_{\Section}\right)\cdot\bar{\Xi}_\kappa\bm{\gamma}\]

But there's a problem: \(c_\vartheta\) is determined by the
\textbf{curvature} \(\varkappa\), not \(\bm{\gamma}\). Under pure
flexure with no applied torque, \(\varkappa = c_\vartheta\) comes from
the SV constraint, which depends on \(\bm{b}_{\Section}\).


Define:
\[
\mathbf{g}_{c} := -\frac{1}{\Jsv}\FrameBasis\cdot\int_{\Section} \bm{r}\times(\nabla\bm{\psi}^T + \PoissonFlexure)\,dA
\]
So \(c_\vartheta = \mathbf{g}_{\vartheta}\cdot\bm{b}_{\Section}\).

Similarly, \(\bar{c}_{\vartheta} = \bar{\mathbf{g}}\cdot\bm{b}_{\Section}\)
where:
\[
\bar{\mathbf{g}} := \FrameBasis\cdot\bar{\bm{I}}^{-1} \int_{\Section} (\bm{r}-\bar{\bm{p}})\times\PoissonFlexure\,dA
\]

Then \(\tilde{\eta}_{\vartheta}\) is linear in \(\bm{\gamma}\), and
\[
\tilde{\eta}_\vartheta = (\mathbf{g}_{\vartheta} - \bar{\mathbf{g}})\cdot\bm{b}_{\Section} = (\mathbf{g}_{\vartheta} - \bar{\mathbf{g}})\cdot\bar{\Xi}_\kappa\bm{\gamma}
\qquad\text{and}\qquad
\mathbf{L}_{\rm wn}^{(\vartheta)} = -(\mathbf{g}_{\vartheta} - \bar{\mathbf{g}})\cdot\bar{\Xi}_\kappa
\]

This is a \(1\times 2\) row vector (scalar constraint from 2D
\(\bm{\gamma}\)).


\subsection{Summary}

\paragraph{V2}

\begin{longtable}[]{@{}
  >{\raggedright\arraybackslash}p{(\linewidth - 8\tabcolsep) * \real{0.1333}}
  >{\raggedright\arraybackslash}p{(\linewidth - 8\tabcolsep) * \real{0.1556}}
  >{\raggedright\arraybackslash}p{(\linewidth - 8\tabcolsep) * \real{0.2222}}
  >{\raggedright\arraybackslash}p{(\linewidth - 8\tabcolsep) * \real{0.2000}}
  >{\raggedright\arraybackslash}p{(\linewidth - 8\tabcolsep) * \real{0.2889}}@{}}
\toprule\noalign{}
\begin{minipage}[b]{\linewidth}\raggedright
Mode
\end{minipage} & \begin{minipage}[b]{\linewidth}\raggedright
Shape
\end{minipage} & \begin{minipage}[b]{\linewidth}\raggedright
Bimoment
\end{minipage} & \begin{minipage}[b]{\linewidth}\raggedright
Bishear
\end{minipage} & \begin{minipage}[b]{\linewidth}\raggedright
Equilibrium
\end{minipage} \\
\midrule\noalign{}
\endhead
\bottomrule\noalign{}
\endlastfoot
\(v\) & \(\mathbf{1}_{\Section}\) & \(\ShearForce\) & \(0\) & \(\ShearForce' = 0\) \\
\(r\) & \(\FrameBasis\otimes\bm{r}\) & \(Ex\,\bm{I}_m\bm{b}_{\Section}\) &
\(\ShearForce\) & \(E\bm{I}_m\bm{b}_{\Section} = \ShearForce\) \\
\end{longtable}

\paragraph{V3}
\begin{longtable}[]{@{}
  >{\raggedright\arraybackslash}p{(\linewidth - 10\tabcolsep) * \real{0.1091}}
  >{\raggedright\arraybackslash}p{(\linewidth - 10\tabcolsep) * \real{0.1273}}
  >{\raggedright\arraybackslash}p{(\linewidth - 10\tabcolsep) * \real{0.2000}}
  >{\raggedright\arraybackslash}p{(\linewidth - 10\tabcolsep) * \real{0.2182}}
  >{\raggedright\arraybackslash}p{(\linewidth - 10\tabcolsep) * \real{0.1818}}
  >{\raggedright\arraybackslash}p{(\linewidth - 10\tabcolsep) * \real{0.1636}}@{}}
\toprule\noalign{}
\begin{minipage}[b]{\linewidth}\raggedright
Mode
\end{minipage} & \begin{minipage}[b]{\linewidth}\raggedright
Shape
\end{minipage} & \begin{minipage}[b]{\linewidth}\raggedright
Amplitude
\end{minipage} & \begin{minipage}[b]{\linewidth}\raggedright
Constraint
\end{minipage} & \begin{minipage}[b]{\linewidth}\raggedright
Bimoment
\end{minipage} & \begin{minipage}[b]{\linewidth}\raggedright
Bishear
\end{minipage} \\
\midrule\noalign{}
\endhead
\bottomrule\noalign{}
\endlastfoot
\(v\) & \(\mathbf{1}_{\Section}\) & \(\alpha'_v = -\bm{\eta}_v\) &
\(-\SectionAverage{\PoissonFlexure}\bar{\Xi}_\kappa\bm{\gamma}\) & \(\ShearForce\) &
\(0\) \\[0.3cm]
\(r\) & \(\FrameBasis\otimes\bm{r}\) & \(\alpha_r = -\bm{\eta}_{\Section}\) &
\(-\mathbf{B}\bar{\Xi}_\kappa\bm{\gamma}\) & \(E\bm{I}_m\bm{b}_{\Section}\)
& \(\ShearForce\) \\[0.3cm]
\(\vartheta\) & \(\FrameBasis\times\bm{r}\) &
\(\alpha'_\vartheta = -\tilde{\eta}_\vartheta\) &
\(-(\mathbf{g}_{\vartheta} - \bar{\mathbf{g}})\bar{\Xi}_\kappa\bm{\gamma}\) & \(T\) &
\(0\) \\
\end{longtable}
with $\mathbf{B} = \mathbf{i} \times \bar{\Theta}\left[\mathbf{i} \psi_\alpha\right] \otimes \mathbf{i}_\alpha$